\chapter{Methodologie}
\label{ch:methodologie}

\section{Onderzoeksopzet}

Deze bachelorproef is opgebouwd volgens de principes van Design Science Research (DSR) \autocite{Peffers2007,Hevner2004}. Deze methodologie is geschikt voor dit onderzoek omdat de focus ligt op het iteratief ontwerpen en technisch realiseren van een oplossing voor de beperkte schaalbaarheid van huidige Basic Life Support (BLS)-training. Hoewel de ontwikkelde proof of concept zich specifiek toespitst op de reanimatieprocedure, is de onderliggende architectuur generiek opgezet zodat toekomstige uitbreiding mogelijk blijft.

Binnen dit onderzoeksproces ligt de nadruk op de fase waarin de oplossing technisch wordt uitgewerkt. Eerst wordt de problematiek van BLS-onderwijs in kaart gebracht via literatuuronderzoek. Vervolgens worden functionele en technische vereisten gedefinieerd op basis van actuele richtlijnen \autocite{ERC2025BLS}. Het zwaartepunt ligt bij de ontwikkeling van de PoC en de technische verificatie van de software-architectuur binnen de afgebakende scenescope.

\section{Onderzoeksstrategie}

Gezien de beperkte tijdlijn voor empirisch onderzoek bij grote groepen gebruikers hanteert deze studie een strategie die gericht is op technische verificatie. In plaats van brede gebruikerstesten wordt de werking van de applicatie getoetst aan de procedurele kaders van de ERC- en AHA-richtlijnen \autocite{ERC2025BLS, AHA2025Guidelines}.

Binnen de effectieve uitvoering van deze bachelorproef werd externe expertreview niet gerealiseerd door tijdsbeperking. De uitgevoerde evaluatie steunt daarom op interne technische verificatie, gestructureerde scenario-walkthroughs en vergelijking van logoutput met vooraf gedefinieerde protocolstappen.

\section{Ontwerp en ontwikkeling van de applicatie}

De XR-toepassing wordt ontwikkeld voor de Meta Quest 3 met Unity als engine \autocite{UnityDocs}. Het ontwerp vertaalt de procedurele stappen van de actuele richtlijnen naar een gelaagde Finite State Machine binnen de software. 

De applicatie voorziet in visuele en auditieve feedback en is technisch zo geconfigureerd dat elke relevante interactie automatisch wordt geregistreerd in een gestructureerd JSON-formaat. De ontwikkeling volgt een iteratief proces waarbij de technische robuustheid en de nauwkeurigheid van de interactie-events centraal staan.

\section{Evaluatie- en validatieopzet}

Vanwege de beperkte scope voor een veldtest wordt de evaluatie van de PoC toegespitst op twee pijlers: procedurele validiteit, met de vraag in hoeverre de geprogrammeerde flow overeenkomt met de officiële reanimatierichtlijnen, en technische performantie, met aandacht voor niet-functionele eisen zoals stabiele runtime en consistente datalogging.

\section{Meetinstrumenten en meetmethoden}

Hoewel grootschalige kwantitatieve dataverzameling buiten de scope valt, zijn de volgende instrumenten gebruikt voor de evaluatie:

De XR-applicatie is voorzien van een geautomatiseerde logging-module. Deze registreert kritieke parameters zoals de volgorde van handelingen, de timing van stappen en de correcte inzet van de AED. De validiteit van deze data wordt gecontroleerd via scenario-walkthroughs door de onderzoeker.

Voor de beoordeling van gebruikerservaring worden de System Usability Scale (SUS) en de Simulator Sickness Questionnaire (SSQ) als theoretisch kader gehanteerd \autocite{SUSBrooke1996, Kennedy1993SSQ}. Deze instrumenten structureren de opzet van een latere gebruikersevaluatie.

Omdat er binnen de beschikbare tijd geen externe reviewronde werd uitgevoerd, worden SUS en SSQ in deze iteratie niet als empirische uitkomstmaten gerapporteerd.

\section{Data-analyse}

De analyse richt zich op de vergelijking tussen feitelijke logdata en voorgeschreven medische protocolstappen. Hierbij wordt onderzocht in hoeverre uitgevoerde handelingen afwijken van de beoogde reanimatiesequentie. Tot slot worden resultaten van technische metingen, zoals runtime-stabiliteit, getoetst aan vooraf gedefinieerde kwaliteitseisen.

\section{Ethische overwegingen}

Hoewel er binnen deze iteratie geen externe expertreviews werden uitgevoerd, blijft het ethisch kader ongewijzigd voor vervolgonderzoek. Bij toekomstige validaties nemen deelnemers vrijwillig deel en wordt feedback geanonimiseerd verwerkt. Daarnaast worden geen patiëntgegevens of gevoelige studentendata verzameld in de huidige PoC-evaluatie.

\section{Beperkingen van de studie}

Binnen de scope van deze bachelorproef is de evaluatie beperkt tot proof-of-conceptvalidatie. Door tijdsbeperkingen is geen uitgebreide gebruikerstest uitgevoerd bij een significante steekproef van studenten. De resultaten tonen technische en procedurele haalbaarheid, maar vormen geen sluitend bewijs voor leerwinst op lange termijn.